% Sections 3.1 to 3.9, from lectures/L03/appendix/a_flip_flops_and_registers.md: one section per
% part, A.1 to A.9.
%
% The four block drawings are drawn with box-drawing characters, so they are textdrawing blocks.

\section{From Combinational to Sequential Logic}
\label{c3:a1}
Every circuit in Chapters~\ref{c1:ch} and~\ref{c2:ch} was \textbf{combinational}: its output is a
pure function of its present inputs, so the same inputs always give the same output, and there is
no notion of ``before'' or ``after'' inside the circuit.

Most useful systems need more. A counter must remember its count. A register must hold a value after
the signal that produced it has changed. An LED toggled by a button must remember whether it is on,
since that is not encoded in the button's momentary state at all.

All of these need \textbf{memory}: an output that depends on the circuit's own \emph{previous}
output. Sequential logic is combinational logic with part of the output fed back into the input, and
that \textbf{feedback loop} is what turns a stateless gate network into a circuit with state. This
chapter builds the primitive that makes controlled feedback practical, the D flip-flop, then the two
things you build directly out of it: registers and edge detectors.

\section{The D Latch}
\label{c3:a2}
The simplest storage element is the \textbf{D latch}: two cross-coupled gates, gated by an
\code{enable} input, storing a single bit.

\needspace{9\baselineskip}
\begin{textdrawing}
        ┌────────────┐
   D ──►│            │──► Q
        │  D LATCH   │
enable─►│            │──► Qn
        └────────────┘
\end{textdrawing}

\noindent\code{D} is the value to store, \code{enable} controls whether the latch is \emph{open}
(transparent) or \emph{closed} (holding), \code{Q} is the stored value and \code{Qn} is always its
inverse.

Each output feeds back into the other's equation, which is exactly \secref{c3:a1}'s feedback loop:
\[ Q = (D' \cdot \bool{enable} + \bool{Qn})' \]
\[ \bool{Qn} = (D \cdot \bool{enable} + Q)' \]

\begin{booktable}{@{}p{3.6em}p{13em}L@{}}
  \thead{\code{enable}} & \thead{Behaviour} & \thead{Result} \\
  \midrule
  \code{1} & latch is \textbf{open} (transparent) & \code{Q = D}, \code{Qn = D'}; output follows
    the input immediately \\
  \code{0} & latch is \textbf{closed} (locked) & \code{Q} and \code{Qn} hold whatever they last
    had \\
\end{booktable}

\noindent\textbf{Build it in CircuitVerse} from one NOT gate, two AND gates and two NOR gates (a NOR
is an OR followed by a NOT). Cross-couple the feedback so the NOR producing \code{Q} feeds the NOR on
the \code{Qn} side and vice versa. Drive \code{D} and \code{enable} from input switches and watch
\code{Q} and \code{Qn} on LEDs.

Then test it:
\begin{itemize}
  \item With \code{enable = 1}, change \code{D} and confirm \code{Q} follows immediately. The latch
    is \textbf{transparent}.
  \item With \code{enable = 0}, confirm \code{Q} is unchanged whatever \code{D} does. The latch is
    \textbf{locked}.
  \item Cycle \code{D} a few times while toggling \code{enable}, and confirm the latch always freezes
    whatever \code{Q} held at the instant \code{enable} fell. That is what
    \textbf{level-sensitive} means: transparent for the entire time \code{enable} is high, not just
    at one instant.
\end{itemize}
That last point is also the latch's fundamental weakness. Because it is transparent for the whole
time \code{enable = 1}, any glitch on \code{D} during that window passes straight through to
\code{Q}. In a large synchronous system, guaranteeing stable data throughout an entire clock phase
is far harder than guaranteeing it around a single instant. That is the problem the D flip-flop
solves.

\section{The D Flip-Flop}
\label{c3:a3}
A \textbf{D flip-flop} stores a bit the same way a latch does, but is \textbf{edge-triggered} rather
than level-sensitive: it samples \code{D} only at the \emph{instant} the clock transitions, not for
the whole time the clock is high or low.

Feed the same \code{D} and the same clock to both and the difference is immediate
(Figure~\ref{c3:fig:latch_vs_flipflop}).

\bookfigure{latch_vs_flipflop}{A latch and a flip-flop fed the same \code{D} and clock: the latch
follows \code{D} throughout every clock-high window, while the flip-flop changes only on the rising
edges.}{c3:fig:latch_vs_flipflop}

\noindent The latch changes three times: twice because \code{D} moved while the clock was high, and
once because the clock went high onto a \code{D} that had already changed. The flip-flop changes
once, because only one of the four rising edges found \code{D} different from the value \code{Q} was
already holding. Both are behaving correctly; they are answering different questions.

\needspace{9\baselineskip}
\begin{textdrawing}
        ┌───────────────┐
   D ──►│               │──► Q
clock ─►│  D FLIP-FLOP  │
enable─►│ (rising-edge) │──► Qn
reset_n►│               │
        └───────────────┘
\end{textdrawing}

\noindent\code{clock} is the timing reference (\secref{c3:a4}). \code{enable} gates whether the
flip-flop accepts a new value \emph{at the triggering edge}. \code{reset_n} is an asynchronous,
active-low reset: whenever it is \code{0}, \code{Q} is forced to \code{0} immediately, regardless of
the clock.

\begin{booktable}{@{}>{\raggedright\arraybackslash}p{17em}L@{}}
  \thead{Condition} & \thead{Result} \\
  \midrule
  \code{reset_n = 0} & \code{Q = 0}, \code{Qn = 1}, immediately (asynchronous, ignores the clock
    entirely) \\
  \code{reset_n = 1}, rising edge on \code{clock}, \code{enable = 1} & \code{Q = D},
    \code{Qn = D'}, sampled at that instant \\
  \code{reset_n = 1}, rising edge on \code{clock}, \code{enable = 0} & \code{Q} holds its previous
    value \\
  \code{reset_n = 1}, no rising edge on \code{clock} & \code{Q} holds its previous value, no matter
    what \code{D} does \\
\end{booktable}

\noindent That last row is worth seeing rather than reading (Figure~\ref{c3:fig:dff_timing}).

\bookfigure{dff_timing}{A flip-flop between edges: \code{D} rises mid-cycle and is captured at the
next rising edge, then dips and recovers entirely between two edges without \code{Q} ever
moving.}{c3:fig:dff_timing}

\noindent\code{D} dips low and comes back inside a single clock period, and \code{Q} never moves,
because no edge happened while it was low. A latch in the same position would have followed it down
and back up.

Internally, an edge-triggered flip-flop is two D latches in series, a \textbf{master-slave} pair. One
is transparent while the clock is low and captures \code{D}; the other is transparent while the clock
is high and passes the captured value to \code{Q}. The net effect is that \code{Q} changes only at
the instant the clock goes low to high, never while it sits at either level. That external behaviour
is what VHDL's \code{rising_edge()} models (\secref{c3:a6}).

\textbf{Build it in CircuitVerse} by extending your \secref{c3:a2} latch:
\begin{itemize}
  \item Duplicate the latch so two sit in series: the \emph{master} feeds the \emph{slave}, whose
    \code{Q} is the flip-flop's output.
  \item Add \code{clock}, \code{reset_n}, \code{D} and \code{enable} as inputs.
  \item Drive the two latches' transparency from the clock and its inverse. To keep the two meanings
    apart, call that internal per-latch control \code{gate}: the master's \code{gate} is
    \code{clock'} and the slave's is \code{clock}. The flip-flop's external \code{enable} is a
    different signal and is wired separately, below.
  \item Wire the external \code{enable} as a 2-to-1 mux (\secref{c2:a4}) on the master's data input:
    when \code{enable = 1} the mux passes \code{D}, and when \code{enable = 0} it passes the
    flip-flop's own \code{Q} back in, so the edge re-captures the value already stored. Gate the
    \emph{data}, never the clock.
  \item Add the reset to \textbf{both} latches. Forcing \code{Q} low takes two changes in each
    latch, not one:
    \begin{itemize}
      \item feed an active-high \code{reset} (\code{reset_n} through a NOT) into the NOR producing
        \code{Q}, forcing \code{Q} to \code{0}.
      \item AND \code{reset_n} into that same latch's \code{D AND gate} term, which frees \code{Qn}
        to go high.
    \end{itemize}
  \item Set the clock period to something slow enough to watch, e.g.\ \code{1000 ms}.
\end{itemize}
With the reset wired in, each latch's equations become:
\[ Q = (D' \cdot \bool{gate} + \bool{Qn} + \bool{reset})' \]
\[ \bool{Qn} = (D \cdot \bool{gate} \cdot \bool{reset\_n} + Q)' \]
A tempting shortcut is to leave the \code{Qn} equation alone and break the loop instead, ANDing
\code{reset_n} onto the \code{Qn} wire feeding \code{Q}'s NOR. That is not enough. With
\code{gate = 1} and \code{D = 1} it leaves \code{Qn} at \code{0}, which is what holds \code{Q} at
\code{1}, so the reset silently does nothing in exactly the case you are most likely to test it in.
Both changes, in both latches, or neither.

Confirm that \code{Q} updates only at the instant the clock rises, holding steady between edges even
if \code{D} changes, and that pulling \code{reset_n} low forces \code{Q} to \code{0} immediately,
with releasing it changing nothing until the next rising edge.

\section{Clocking}
\label{c3:a4}
The \textbf{clock} is a periodic square wave, read by every flip-flop in a synchronous circuit from
the same source. Two numbers describe it: the \textbf{period} \code{T}, the time for one full cycle,
and the \textbf{frequency} \code{f = 1/T} in Hz.

Each period has exactly two \textbf{edges}: a \textbf{rising edge} from \code{0} to \code{1}, and a
\textbf{falling edge} from \code{1} to \code{0}. A synchronous design picks one of them, almost
always the rising edge, and updates every flip-flop on that edge and only that edge. Using a single
edge everywhere is what makes a large circuit's timing analyzable: every flip-flop changes at the
same well-defined instants, so you can reason one clock cycle at a time rather than tracking every
signal continuously.

Simulated and real clocks run on wildly different timescales, and the logic is identical either way:
\begin{itemize}
  \item In CircuitVerse, set a period you can watch, e.g.\ \code{1000 ms}.
  \item On the DE0-CV the system clock is a fixed \code{50 MHz}, a period of \code{20 ns}. Nobody
    perceives individual edges at that rate, so you reason in terms of ``on the next rising edge'',
    never in wall-clock time.
\end{itemize}

\section{Registers: Flip-Flops in Parallel}
\label{c3:a5}
A \textbf{register} stores more than one bit: N D flip-flops side by side, one per bit, all sharing
the same \code{clock}, \code{reset_n} and \code{enable}.

\needspace{9\baselineskip}
\begin{textdrawing}
D(0) ──►[D FF]──► Q(0)
D(1) ──►[D FF]──► Q(1)
D(2) ──►[D FF]──► Q(2)
   ⋮        ⋮
D(N-1)──►[D FF]──► Q(N-1)

  (clock, reset_n, enable shared by every flip-flop above)
\end{textdrawing}

\noindent That is the entire idea: an 8-bit register is eight D flip-flops updating on the same clock
edge.

In VHDL you rarely instantiate N flip-flops explicitly. Declare one \code{std_logic_vector} and
assign the whole thing inside a single synchronous process (\secref{c3:a6}): every bit synthesizes
to its own flip-flop, all sharing the process's clock and reset.

\needspace{14\baselineskip}
\begin{vhdlcode}
signal q: std_logic_vector(3 downto 0);
...
process (clock, reset_n) is
begin
    if (reset_n = '0') then
        q <= "0000";
    elsif (rising_edge(clock)) then
        if (enable = '1') then
            q <= d; -- d is a std_logic_vector(3 downto 0) input
        end if;
    end if;
end process;
\end{vhdlcode}

\noindent Four bits, one process, four flip-flops in hardware: no different in principle from four
\code{d_flip_flop} instances wired in parallel.

\section{The Synchronous Process Template in VHDL}
\label{c3:a6}
Every synchronous element in this chapter, the flip-flop, the register and the edge detector, uses
the same shape. Without a reset, a single D flip-flop is:

\begin{vhdlcode}
process (clock) is
begin
    if (rising_edge(clock)) then
        q <= d;
    end if;
end process;
\end{vhdlcode}

\noindent\code{q} takes \code{d}'s value at every rising edge and holds it the rest of the time. That
one line inside the \code{if} is enough for synthesis to infer a real D flip-flop.

\subsection{Rules}
\label{c3:rules}
\begin{itemize}
  \item \textbf{\code{clock} is the only signal in the sensitivity list}, unless there is an
    asynchronous reset. The process should run only when the clock changes; it has no reason to
    react to anything else.
  \item A process reacts to \emph{any} change on a sensitivity-list signal, not just rising ones, so
    \code{rising_edge()} filters for the rising edge specifically. All logic that updates
    synchronously goes inside that \code{if}.
  \item Every signal assigned inside such a process synthesizes to a flip-flop, one bit each.
\end{itemize}

\subsection{With an asynchronous reset}
\label{c3:asyncreset}
An asynchronous reset must also appear in the sensitivity list, since it can be asserted
independently of the clock, and it takes priority when active.

\needspace{10\baselineskip}
\begin{vhdlcode}
process (clock, reset_n) is
begin
    if (reset_n = '0') then
        q <= '0';
    elsif (rising_edge(clock)) then
        q <= d;
    end if;
end process;
\end{vhdlcode}

\noindent Read it in priority order: check the reset first, and only if it is inactive check for a
rising edge. If neither holds (a falling edge, or \code{reset_n} going from \code{0} back to
\code{1}, which also re-triggers the process), nothing is assigned and \code{q} keeps its previous
value.

\subsection{When the assignment actually happens}
\label{c3:scheduling}
One rule sits underneath all of the above, and it is the single place where \code{<=} behaves unlike
assignment in any language you have written software in.

A signal assignment does not take effect when its line runs. It \textbf{schedules} a value, applied
only once the process finishes this pass and suspends. Until then, every read of that signal returns
the value it held \emph{before the pass started}, however many assignments have already executed.

That sounds like a technicality. It is the reason the template above is a flip-flop rather than a
wire:

\begin{vhdlcode}
process (clock) is
begin
    if (rising_edge(clock)) then
        b <= a;
        c <= b;
    end if;
end process;
\end{vhdlcode}

\noindent Read as software, \code{b} and \code{c} both end up holding \code{a}. In hardware they do
not. On each edge \code{b} is scheduled to take \code{a}'s value, and \code{c} is scheduled to take
the value \code{b} held \emph{coming into this edge}, not the one just scheduled onto it. The result
is two flip-flops in a chain: \code{a} reaches \code{b} after one edge and \code{c} after two. Swap
the lines and nothing changes, because neither assignment can observe the other's result.

Two consequences worth carrying forward:
\begin{itemize}
  \item \textbf{Order does not matter} between assignments to \emph{different} signals in the same
    clocked process. They all read the old values and all take effect together. That is the opposite
    of the sequential reading the word ``process'' invites, and it is why a chain of flip-flops can
    be written in any order.
  \item \textbf{Assign the same signal twice in one pass and only the last assignment survives.} The
    earlier ones are discarded without ever reaching the signal.
\end{itemize}
\Secref{c3:a7}'s edge detector is exactly the two-element chain above, and works for exactly this
reason. Chapter~\ref{c5:ch} returns to this rule from the other side, with a \code{variable}, which
updates immediately and so behaves the way the software reading expects.

\section{Edge Detection}
\label{c3:a7}
A flip-flop's other everyday job is detecting the \emph{instant} a signal changes rather than reading
its level: store the signal's previous value in a flip-flop, then compare it against the current
value with a gate.
\begin{itemize}
  \item \textbf{Rising-edge detection} (\code{0} to \code{1}):
    \[ \bool{edge} = \bool{current} \cdot \bool{previous}' \]
  \item \textbf{Falling-edge detection} (\code{1} to \code{0}):
    \[ \bool{edge} = \bool{current}' \cdot \bool{previous} \]
\end{itemize}
The rising-edge form as a circuit; swap which side is inverted for a falling edge, which is the form
\secref{c3:a9}'s active-low button needs:

\needspace{9\baselineskip}
\begin{textdrawing}
signal ──┬───────────────────────────► AND ──► edge
         │                              ▲
         └──►[D FF]── previous ──(NOT)──┘
                ▲
           clock, reset_n
\end{textdrawing}

\noindent\code{edge} is high for exactly one clock cycle, the one immediately after the transition,
because \code{previous} only catches up on the \emph{next} edge. This is the mechanism behind ``do
something once per button press'' rather than ``do something every clock cycle the button is held'',
which is exactly what \secref{c3:a9} solves.

\section{Worked Example: A D Flip-Flop in VHDL}
\label{c3:a8}
\file{lectures/L03/d_flip_flop/d_flip_flop.vhd} implements \secref{c3:a3} and \secref{c3:a6}
combined: an enable gating whether a rising edge captures \code{d}, and an asynchronous active-low
reset.

\begin{vhdlcode}
entity d_flip_flop is
    port(clock, reset_n, d, enable: in std_logic;
         q, q_n                   : out std_logic);
end entity;
\end{vhdlcode}

\noindent Its single process is the reset-plus-enable variant of \secref{c3:a6}'s template: reset
takes priority, then a rising edge captures \code{d} into an internal signal \code{q_s} if
\code{enable = '1'}. \code{q_s} exists because an \emph{output} port cannot be read back inside the
same architecture; \code{q} and \code{q_n} are both driven combinationally from it, outside the
process. This is the whole of \file{d_flip_flop.vhd}:

\vhdlfile{lectures/L03/d_flip_flop/d_flip_flop.vhd}

To try it on hardware, assign \code{clock} to the \code{50 MHz} system clock, \code{reset_n} to a
pushbutton, \code{d} and \code{enable} to two slide switches, and \code{q}/\code{q_n} to two LEDs.
Then flip \code{enable} on and change \code{d}, and the LED follows; flip \code{enable} off and change
\code{d}, and it does not.

There is no clock edge you can perceive at \code{50 MHz}, so what you are really confirming is that
the logic still holds once you can no longer see the individual edges. The flip-flop looks like it
responds to your switch; in fact it responded to one of fifty million clock edges that second, and
the switch only decided which value that edge captured.

\section{Worked Example: Edge-Detected LED Toggle}
\label{c3:a9}
\file{lectures/L03/led_toggle/led_toggle.vhd} combines \secref{c3:a5}, \secref{c3:a6} and
\secref{c3:a7} into one small system: two independent push buttons, each toggling its own LED
exactly once per press.

\begin{vhdlcode}
entity led_toggle is
    port(clock, reset_n: in std_logic;
         button_n      : in std_logic_vector(1 downto 0);
         led           : out std_logic_vector(1 downto 0));
end entity;
\end{vhdlcode}

\noindent Three 2-bit signals, two held in clocked processes and one computed between them:
\begin{itemize}
  \item \code{button_prev_n} is a 2-bit register (\secref{c3:a5}) holding last cycle's
    \code{button_n}, reset to \code{"11"} since the buttons are active-low.
  \item \code{button_edge <= (not button_n) and button_prev_n} is \secref{c3:a7}'s falling-edge
    detector applied to both bits at once, high for one cycle exactly when a button goes from
    released to pressed.
  \item \code{led_s} is a second 2-bit register which, instead of capturing a new value each cycle,
    \textbf{toggles} bit \code{i} only when \code{button_edge(i) = '1'}. Without that gating it would
    toggle every clock cycle instead of once per press, which is the whole point of computing
    \code{button_edge}.
\end{itemize}
This is the whole of \file{led_toggle.vhd}:

\vhdlfile{lectures/L03/led_toggle/led_toggle.vhd}

\textbf{On the board}, as the course demonstrates it: create the Quartus Prime Lite project targeting
the DE0-CV's \code{5CEBA4F23C7N} exactly as in Chapter~\ref{c1:ch} (Appendix~\ref{app:quartus}),
then assign \code{clock} to the \code{50 MHz} oscillator, \code{reset_n} and
\code{button_n(1 downto 0)} to three push buttons, and \code{led(1 downto 0)} to two LEDs. Each
press should toggle exactly one LED, once, no matter how long you hold it.

\textbf{A caveat, deliberately left open:} \code{button_n} is a raw signal straight off a physical
pushbutton pin. In simulation that is a clean, instantaneous transition; on real hardware it is
neither instantaneous (contacts bounce) nor synchronized to \code{clock}. This edge detector is
logically correct but \emph{unsafe} to wire directly to a real button on an FPGA without more work
first, and that work is exactly what Chapter~\ref{c4:ch} covers.
